\section{The level one measure}\label{lo:sec:level}

Throughout, $\xi(s)=\tfrac12s(s-1)\pi^{-s/2}\Gamma(s/2)\zeta(s)$, zeros are $\rho=\sigma_\rho+i\gamma_\rho$ counted with multiplicity, and $\theta_0=\sup\sigma_\rho\le1$. For every $\theta\ge\theta_0$, that is for every $\theta$ such that $\zeta$ has no zeros in $\Re s>\theta$ (in particular for every $\theta\ge1$, unconditionally), let
\begin{equation}\label{lo:eq:Utheta}
\mathfrak U_\theta=\sum_{\rho}\mathrm{Cauchy}(\gamma_\rho,\theta-\sigma_\rho),\qquad\mathrm{Cauchy}(\gamma,0)=\delta_\gamma,
\end{equation}
the sum over the zeros of both half planes, a positive measure on $\R$ in the variable $v$ with $\int\mathfrak U_\theta(dv)/(1+v^2)<\infty$. It is twice the symmetric measure $\widetilde U_\theta$ of Theorem~\ref{lv:thm:semigroup}, and the image of its restriction to $v>0$ under $v\mapsto v^2$ is the Thorin measure $U_\theta$, each conjugate pair of zeros contributing one term.

\begin{theorem}[density of the Thorin measure; at $\theta=1$, $v\ne0$, Propositions~\ref{bp:prop:thorin} and~\ref{bp:prop:sin2}]\label{lo:thm:density}
If $\theta>\sigma_\rho$ for every $\rho$ (for instance if $\theta>\theta_0$, and for every $\theta\ge1$ because $\zeta(1+it)\ne0$), the measure $\mathfrak U_\theta$ is absolutely continuous with real analytic density
\begin{equation}\label{lo:eq:dens}
u_\theta(v)=\frac1\pi\Re\frac{\xi'}{\xi}(\theta+iv)=\frac1\pi\,\partial_\sigma\log|\xi(\sigma+iv)|\Big|_{\sigma=\theta},
\end{equation}
and $\mathfrak U_\theta((0,V])=\frac1\pi\arg\xi(\theta+iV)$ for the branch with $\arg\xi(\theta)=0$. At $\theta=1$,
\begin{equation}\label{lo:eq:u1}
\begin{split}u_1(v)&=\frac1\pi\Bigl[\frac1{1+v^2}-\frac{\log\pi}2+\frac12\Re\frac{\Gamma'}{\Gamma}\Bigl(\frac{1+iv}2\Bigr)\\&\qquad-\lim_{X\to\infty}\Bigl(\sum_{n\le X}\frac{\Lambda(n)}{n}\cos(v\log n)-\frac{\sin(v\log X)}{v}\Bigr)\Bigr]\end{split}
\end{equation}
for every real $v$, where $\sin(v\log X)/v$ means $\log X$ at $v=0$. The limit exists, locally uniformly in $v$, and equals $\Re C(1+iv)$ with $C(s)=-\zeta'(s)/\zeta(s)-1/(s-1)$: it is $-\Re(\zeta'/\zeta)(1+iv)$ for $v\ne0$ and minus Euler's constant at $v=0$, while the cosine series itself has no limit for any real $v$ (compare Proposition~\ref{bp:prop:sin2}). For $\theta>1$,
\[
\begin{split}u_\theta(v)&=\frac1\pi\Bigl[\frac\theta{\theta^2+v^2}+\frac{\theta-1}{(\theta-1)^2+v^2}-\frac{\log\pi}2\\&\qquad+\frac12\Re\frac{\Gamma'}{\Gamma}\Bigl(\frac{\theta+iv}2\Bigr)-\sum_{n\ge2}\frac{\Lambda(n)}{n^\theta}\cos(v\log n)\Bigr],\end{split}
\]
the prime sum converging absolutely.
\end{theorem}

\begin{proof}
Each term of \eqref{lo:eq:Utheta} is a probability measure, and for $0\le c\le\theta$
\[
\int\frac{P_c(v-\gamma)}{1+v^2}\,dv=\pi P_{1+c}(\gamma)=\frac{1+c}{(1+c)^2+\gamma^2}\le\frac{1+\theta}{1+\gamma^2};
\]
since $\sum_\rho(1+\gamma_\rho^2)^{-1}<\infty$, the sum converges to a positive measure with $\int\mathfrak U_\theta(dv)/(1+v^2)<\infty$. If every width $\theta-\sigma_\rho$ is positive, fix $R>0$. The finitely many zeros with $|\gamma_\rho|\le R+2\theta+2$ give Cauchy densities of positive width, which are real analytic; for the others the densities, continued to $|\Re z|\le R$, $|\Im z|\le\tfrac12$, are holomorphic and bounded by $2\theta/(\pi(|\gamma_\rho|-R)^2)$, so their sum is holomorphic there. Hence $\mathfrak U_\theta$ has a real analytic density.

By the Hadamard product with conjugate pairing, $(\xi'/\xi)(s)=\sum_\rho1/(s-\rho)$, so
\[
\frac1\pi\Re\frac{\xi'}{\xi}(\theta+iv)=\frac1\pi\sum_\rho\frac{\theta-\sigma_\rho}{(\theta-\sigma_\rho)^2+(v-\gamma_\rho)^2},
\]
which is the density of \eqref{lo:eq:Utheta}. The second form in \eqref{lo:eq:dens} is the Cauchy Riemann equation for $\log\xi$, and integrating $\partial_v\arg\xi=\partial_\sigma\log|\xi|$ gives the mass formula. For \eqref{lo:eq:u1}, $\xi'/\xi=1/s+1/(s-1)-\tfrac12\log\pi+\tfrac12\Gamma'/\Gamma(s/2)+\zeta'/\zeta$, so $u_1=g_1-\pi^{-1}\Re C(1+iv)$, where $g_1$ collects the first three terms of \eqref{lo:eq:u1} and $C(s)=-\zeta'(s)/\zeta(s)-1/(s-1)$ is analytic on a neighbourhood of $\Re s\ge1$, because $\zeta(1+it)\ne0$ and the pole of $\zeta$ cancels; at $v=0$ this is the value of the continuous function $u_1$. Write $\sum_{n\le t}\Lambda(n)=t+R(t)$ with $R(t)=O(te^{-c\sqrt{\log t}})$. Partial summation gives, for $\Re s\ge1$,
\[
\sum_{n\le X}\frac{\Lambda(n)}{n^s}-\int_1^X\frac{dt}{t^s}=\frac{R(X)}{X^s}-R(1)+s\int_1^X\frac{R(t)}{t^{s+1}}\,dt ,
\]
and as $X\to\infty$ the right side converges, uniformly for $s$ in compact subsets of $\Re s\ge1$, because $\int^\infty e^{-c\sqrt{\log t}}\,dt/t<\infty$. For $\Re s>1$ the left side tends to $C(s)$, so by continuity the limit is $C(s)$ on $\Re s=1$ too. At $s=1+iv$ the real part of $\int_1^Xt^{-s}\,dt$ is $\sin(v\log X)/v$, and $\log X$ at $v=0$, which has no limit as $X\to\infty$; this gives the compensated limit in \eqref{lo:eq:u1} and the divergence of the cosine series. For $\theta>1$ the Dirichlet series converges absolutely, and $\Re1/(s-1)=(\theta-1)/((\theta-1)^2+v^2)$ at $s=\theta+iv$.
\end{proof}

At $v=0$ the compensated limit is $\lim(\sum_{n\le X}\Lambda(n)/n-\log X)=-\gamma_E$, and \eqref{lo:eq:u1} gives $u_1(0)=(1+\gamma_E/2-\log2-\tfrac12\log\pi)/\pi=(2+\gamma_E-\log4\pi)/(2\pi)=0.00735159249234$, the value of the continuous function $u_1$ at the origin (computed, to $30$ digits; the values $u_1(10^{-3})-u_1(0)=3.5\times10^{-11}$ and $u_1(0.05)-u_1(0)=8.8\times10^{-8}$ confirm it). It is also $(2\pi)^{-1}\sum_\rho1/(\rho(1-\rho))$.

So the level one Thorin measure of the reciprocal xi function is an explicit function of the primes, with no reference to zeros, and a positive measure by Theorem~\ref{lv:thm:thorin}. Positivity of \eqref{lo:eq:u1} is itself a statement about primes: it says that the limit in \eqref{lo:eq:u1} is at most $1/(1+v^2)-\tfrac12\log\pi+\tfrac12\Re\Gamma'/\Gamma((1+iv)/2)$, the right side being $\tfrac12\log(v/2\pi)+O(1/v^2)$. It follows from the Euler product (Proposition~\ref{bp:prop:thorin}). Numerically, the mass of $u_1$ on $[0,400]$ is $201.67$ and there are $202$ zeros below height $400$ (computed: Simpson on step $0.05$ gives $201.66682$, and the formula below over the first $1700$ ordinates plus the smooth tail gives $201.66681$), the difference coming from the Cauchy smoothing of the zero count,
\[
\mathfrak U_1((0,V])=\frac1\pi\sum_{\gamma_\rho>0}\Bigl[\arctan\frac{V-\gamma_\rho}{1-\sigma_\rho}+\arctan\frac{V+\gamma_\rho}{1-\sigma_\rho}\Bigr].
\]

\section{The Fourier side: primes}\label{lo:sec:fourier}

Write $\widehat\mu(\omega)=\int e^{-i\omega v}\mu(dv)$ for the Fourier transform of a tempered measure.

\begin{proposition}[spectrum of the level one measure]\label{lo:prop:spectrum}
For $\theta\ge1$ let $g_\theta(v)=\frac1\pi\bigl[\theta/(\theta^2+v^2)-\tfrac12\log\pi+\tfrac12\Re\frac{\Gamma'}{\Gamma}((\theta+iv)/2)\bigr]$. As tempered distributions,
\begin{equation}\label{lo:eq:spectrum}
\widehat{\mathfrak U_1}(\omega)=\widehat{g_1}(\omega)+1-\sum_{n\ge2}\frac{\Lambda(n)}{n}\bigl(\delta_{\log n}+\delta_{-\log n}\bigr)(\omega),
\end{equation}
the constant $1$ being the transform of the atom $\delta_0$ that the pole of $\zeta$ at $s=1$ contributes in the limit $\theta\downarrow1$; on the zero side $\widehat{\mathfrak U_1}(\omega)=\sum_\rho e^{-(1-\sigma_\rho)|\omega|}e^{-i\omega\gamma_\rho}$, the limit of the sums over $|\gamma_\rho|<T$. Equivalently, for every $\varphi\in\mathcal S(\R)$,
\begin{equation}\label{lo:eq:weil}
\sum_\rho\bigl(P_{1-\sigma_\rho}*\varphi\bigr)(\gamma_\rho)=\int g_1\varphi\,dv+\varphi(0)-\frac1{2\pi}\sum_{n\ge2}\frac{\Lambda(n)}{n}\bigl(\widehat\varphi(\log n)+\widehat\varphi(-\log n)\bigr),
\end{equation}
with $P_c$ the Cauchy kernel of width $c$. For $\theta\ge1$ the semigroup $\mathfrak U_\theta=P_{\theta-1}*\mathfrak U_1$ reads, since $g_\theta=P_{\theta-1}*g_1$, $\widehat{\mathfrak U_\theta}(\omega)=\widehat{P_{\theta-1}*\mathfrak U_1}(\omega)=e^{-(\theta-1)|\omega|}\widehat{\mathfrak U_1}(\omega)=\widehat{g_\theta}(\omega)+e^{-(\theta-1)|\omega|}-\sum_{n\ge2}\Lambda(n)n^{-\theta}(\delta_{\log n}+\delta_{-\log n})(\omega)$, which on the prime side is the replacement of $\Lambda(n)/n$ by $\Lambda(n)n^{-\theta}$.
\end{proposition}

\begin{proof}
For $1<\theta\le2$ the formula of Theorem~\ref{lo:thm:density} reads
\[
u_\theta=g_\theta+P_{\theta-1}-\pi^{-1}\sum_n\Lambda(n)n^{-\theta}\cos(v\log n),
\]
absolutely convergent. Integrating against $\varphi\in\mathcal S(\R)$ and using
\[
\int\varphi(v)\cos(v\log n)\,dv=\tfrac12\bigl(\widehat\varphi(\log n)+\widehat\varphi(-\log n)\bigr),
\]
we get
\[
\begin{split}\sum_\rho(P_{\theta-\sigma_\rho}*\varphi)(\gamma_\rho)&=\int g_\theta\varphi\,dv+(P_{\theta-1}*\varphi)(0)\\&\quad-\frac1{2\pi}\sum_{n\ge2}\frac{\Lambda(n)}{n^\theta}\bigl(\widehat\varphi(\log n)+\widehat\varphi(-\log n)\bigr).\end{split}
\]
Let $\theta\downarrow1$. On the left, splitting the integral at $|v-\gamma|=|\gamma|/2$ gives $|(P_c*\varphi)(\gamma)|\le C_\varphi/(1+\gamma^2)$ uniformly for $0\le c\le2$; each term is continuous in $c$ and $\sum_\rho(1+\gamma_\rho^2)^{-1}<\infty$, so dominated convergence over the zeros applies. On the right, $g_\theta\to g_1$ pointwise with $|g_\theta(v)|\le C\log(2+|v|)$ for $1\le\theta\le2$, $(P_{\theta-1}*\varphi)(0)\to\varphi(0)$, and $\sum_n\Lambda(n)n^{-1}|\widehat\varphi(\pm\log n)|<\infty$ because $\widehat\varphi(\omega)=O((1+|\omega|)^{-3})$ and $\sum_{n\le x}\Lambda(n)/n=\log x+O(1)$. This proves \eqref{lo:eq:weil}. Its left side is $\int\varphi\,d\mathfrak U_1$, and $\int f\varphi\,dv=\frac1{2\pi}\int\widehat f(\omega)\widehat\varphi(-\omega)\,d\omega$, so \eqref{lo:eq:weil} for all $\varphi\in\mathcal S(\R)$ is \eqref{lo:eq:spectrum}. The zero side is the transform of \eqref{lo:eq:Utheta} term by term, $\widehat{\mathrm{Cauchy}(\gamma,c)}(\omega)=e^{-c|\omega|}e^{-i\omega\gamma}$, the sums over $|\gamma_\rho|<T$ converging in $\mathcal S'$ by the same bound. For $\theta>1$ the computation without the limit gives the last formula, and $\mathfrak U_\theta=P_{\theta-1}*\mathfrak U_1$ holds term by term because $P_{\theta-1}*P_c=P_{\theta-1+c}$.
\end{proof}

The identity \eqref{lo:eq:weil} is an explicit formula of Weil's type \cite{Weil52}. For a conjugate pair of zeros $\rho,\bar\rho$, with
\[
H_\varphi(z)=\frac1{2\pi}\int_0^\infty\bigl(\widehat\varphi(t)+\widehat\varphi(-t)\bigr)e^{(z-1)t}\,dt ,
\]
one has $(P_{1-\sigma}*\varphi)(\gamma)+(P_{1-\sigma}*\varphi)(-\gamma)=H_\varphi(\rho)+H_\varphi(\bar\rho)$, so the left side of \eqref{lo:eq:weil} is $\sum_\rho H_\varphi(\rho)$, summed over conjugate pairs: one test function evaluated at the zeros. Under RH it is the classical explicit formula tested against $P_{1/2}*\varphi$ (compare Theorem~\ref{dc:thm:weil}). With Theorem~\ref{dc:thm:weil}, it identifies the Thorin measure with the object of the explicit formula. As a check, for $\varphi(v)=e^{-v^2/2}$ the two sides of \eqref{lo:eq:weil} agree to thirteen digits in the form $\int(u_1-g_1)\varphi\,dv=\varphi(0)-(2\pi)^{-1}\sum_n\Lambda(n)n^{-1}(\widehat\varphi(\log n)+\widehat\varphi(-\log n))$, both equal to $0.4267700651131$ (computed in the review of the source with $120$ point Gauss and Legendre quadrature and $20$ digit arithmetic, not rerun here).

The spectral peaks of $u_1$ sit at $\log2,\log3,\log4,\log5,\log7,\dots$ with heights roughly proportional to $\Lambda(n)/n$. Table~\ref{lo:tab:spectrum} shows them recovered from a Hann windowed Fourier transform of $u_1$ sampled on $[-400,400]$ at spacing $0.05$ (by symmetry), padded to a frequency step of $3\times10^{-5}$.

\begin{table}[ht]
\centering
\small
\begin{tabular}{ccccc}
\toprule
$n$ & peak $\omega$ & $e^{\omega}$ & relative height & $\Lambda(n)/n$ relative to $n=3$\\
\midrule
2 & 0.6931 & 2.000 & 0.950 & 0.946\\
3 & 1.0986 & 3.000 & 1.000 & 1.000\\
4 & 1.3863 & 4.000 & 0.449 & 0.473\\
5 & 1.6095 & 5.000 & 0.872 & 0.879\\
7 & 1.9459 & 7.000 & 0.747 & 0.759\\
8 & 2.0794 & 8.000 & 0.202 & 0.237\\
9 & 2.1972 & 9.000 & 0.303 & 0.333\\
11 & 2.3979 & 11.000 & 0.576 & 0.595\\
13 & 2.5650 & 13.000 & 0.517 & 0.539\\
17 & 2.8332 & 17.000 & 0.430 & 0.455\\
19 & 2.9444 & 19.000 & 0.397 & 0.423\\
23 & 3.1355 & 23.000 & 0.343 & 0.372\\
29 & 3.3673 & 29.000 & 0.286 & 0.317\\
31 & 3.4340 & 31.001 & 0.271 & 0.302\\
\bottomrule
\end{tabular}
\caption{Spectral peaks of the level one Thorin measure (computed). Every peak lies within $10^{-4}$ of the logarithm of a prime or prime power. The heights, relative to that of $n=3$, depend on the window and follow the relative value of $\Lambda(n)/n$ to within $0.04$.}\label{lo:tab:spectrum}
\end{table}

\section{The backward flow}\label{lo:sec:backward}

Fix $\tfrac12\le\theta'<1$, put $a=1-\theta'$, and let $\Delta_\Omega(\omega)=(1-|\omega|/\Omega)_+$ be the Fej\'er window. Let $K_{\theta',\Omega}$ be the inverse Fourier transform of $\Delta_\Omega(\omega)e^{a|\omega|}$; this multiplier is even, continuous, compactly supported and has a derivative of bounded variation, so $K_{\theta',\Omega}$ is smooth and $|K_{\theta',\Omega}(v)|\le C_{\Omega,a}/(1+v^2)$. The band limited backward flow from level one to level $\theta'$ is
\begin{equation}\label{lo:eq:backward}
B_{\theta',\Omega}\mathfrak U_1=\mathfrak U_1*K_{\theta',\Omega},
\end{equation}
a smooth function of growth $O(\log(2+|v|))$, since $\mathfrak U_1([t,t+1])=O(\log(2+|t|))$; formally its transform is $\Delta_\Omega(\omega)e^{a|\omega|}\widehat{\mathfrak U_1}(\omega)$. Let $\Phi_\Omega(v)=(1-\cos\Omega v)/(\pi\Omega v^2)$ be the Fej\'er kernel, the inverse transform of $\Delta_\Omega$.

\begin{theorem}[dichotomy for the backward flow]\label{lo:thm:dichotomy}
Let $\tfrac12\le\theta'<1$ and $a=1-\theta'$.
\begin{enumerate}
\item[(i)] If $\zeta$ has no zeros in $\Re s>\theta'$, then $B_{\theta',\Omega}\mathfrak U_1=\mathfrak U_{\theta'}*\Phi_\Omega$, with $\mathfrak U_{\theta'}$ as in \eqref{lo:eq:Utheta}, which has atoms at the zeros with $\sigma_\rho=\theta'$. As $\Omega\to\infty$, $B_{\theta',\Omega}\mathfrak U_1\to\mathfrak U_{\theta'}$ vaguely, and $\int B_{\theta',\Omega}\mathfrak U_1(v)\,dv/(v^2+s)\to\int\mathfrak U_{\theta'}(dv)/(v^2+s)$ for every $s>0$. If $\theta'>\sigma_\rho$ for every $\rho$, the convergence is locally uniform as densities.
\item[(ii)] If $\zeta$ has a zero $\rho$ with $\sigma_\rho>\theta'$, there is no signed measure $\mu$ with $\int|\mu|(dv)/(1+v^2)<\infty$ such that $\int B_{\theta',\Omega}\mathfrak U_1(v)\,dv/(v^2+s)\to\int\mu(dv)/(v^2+s)$ as $\Omega\to\infty$ for all sufficiently large $s$.
\end{enumerate}
\end{theorem}

\begin{proof}
Write $f_s(v)=1/(v^2+s)$, so that $f_s=(\pi/\sqrt s)P_{\sqrt s}$ and $\widehat{f_s}(\omega)=(\pi/\sqrt s)e^{-\sqrt s|\omega|}$. The measures below are positive or satisfy $\int|\mu|/(1+v^2)<\infty$, and the kernels are $O(1/(1+v^2))$, so every use of Fubini's theorem is absolute.

(i) Every zero has $\sigma_\rho\le\theta'$, so $\mathfrak U_{\theta'}$ is defined and, term by term, $\mathfrak U_1=\mathfrak U_{\theta'}*P_a$, because $P_c*P_a=P_{c+a}$ for $c\ge0$ (Theorem~\ref{lv:thm:semigroup}). The integrable function $P_a*K_{\theta',\Omega}$ has transform $e^{-a|\omega|}\Delta_\Omega(\omega)e^{a|\omega|}=\Delta_\Omega(\omega)$, so $P_a*K_{\theta',\Omega}=\Phi_\Omega$ and $B_{\theta',\Omega}\mathfrak U_1=\mathfrak U_{\theta'}*P_a*K_{\theta',\Omega}=\mathfrak U_{\theta'}*\Phi_\Omega$. The kernel $\Phi_\Omega$ is nonnegative, has integral one and satisfies $\Phi_\Omega(u)\le2/(\pi\Omega u^2)$. For $f$ continuous with support in $[-R,R]$, $\Phi_\Omega*f\to f$ uniformly and $|\Phi_\Omega*f(v)|\le8\|f\|_1/(\pi\Omega v^2)$ for $|v|\ge2R$, so $|\Phi_\Omega*f-f|\le\eta_\Omega/(1+v^2)$ with $\eta_\Omega\to0$, and $\int f\,d(\mathfrak U_{\theta'}*\Phi_\Omega)=\int(\Phi_\Omega*f)\,d\mathfrak U_{\theta'}\to\int f\,d\mathfrak U_{\theta'}$. For $f=f_s$, $\Phi_\Omega*f_s\to f_s$ pointwise and $\Phi_\Omega*f_s\le C_s/(1+v^2)$ for $\Omega\ge1$, and dominated convergence gives the second statement. If $\theta'>\sigma_\rho$ for every $\rho$, $\mathfrak U_{\theta'}$ has a continuous density $u_{\theta'}$ (Theorem~\ref{lo:thm:density}). For $|v|\le R$ and $0<\delta\le1$, $|\Phi_\Omega*u_{\theta'}(v)-u_{\theta'}(v)|$ is at most the largest change of $u_{\theta'}$ over distances $\delta$ on $[-R-1,R+1]$, plus $(2/(\pi\Omega))\int_{|t|\ge\delta}u_{\theta'}(v-t)t^{-2}\,dt+(4/(\pi\Omega\delta))\max_{[-R,R]}u_{\theta'}$, and the last two terms are $O_{R,\delta}(1/\Omega)$.

(ii) Let $s>a^2$ and $b=\sqrt s-a>0$. The integrable function $K_{\theta',\Omega}*f_s$ has transform $(\pi/\sqrt s)\Delta_\Omega(\omega)e^{-b|\omega|}$, so $K_{\theta',\Omega}*f_s=(\pi/\sqrt s)\Phi_\Omega*P_b=(b/\sqrt s)\Phi_\Omega*f_{b^2}$ and
\[
\begin{split}\int B_{\theta',\Omega}\mathfrak U_1(v)f_s(v)\,dv&=\int(K_{\theta',\Omega}*f_s)\,d\mathfrak U_1=\frac b{\sqrt s}\int(\Phi_\Omega*f_{b^2})\,d\mathfrak U_1\\&\longrightarrow\frac b{\sqrt s}\int\frac{\mathfrak U_1(dv)}{v^2+b^2}\end{split}
\]
as $\Omega\to\infty$, by the domination used in (i). Since $\int P_c(v-\gamma)f_{b^2}(v)\,dv=(\pi/b)P_{b+c}(\gamma)=b^{-1}\Re(b+c-i\gamma)^{-1}$, summing over the zeros with $c=1-\sigma_\rho$ gives
\[
\int\frac{\mathfrak U_1(dv)}{v^2+b^2}=\frac1b\sum_\rho\Re\frac1{1+b-\rho}=\frac1b\,\frac{\xi'}{\xi}(1+b),
\]
so the limit is $(\xi'/\xi)(\theta'+\sqrt s)/\sqrt s$. If a measure $\mu$ as in (ii) existed, then for all large $s>0$
\[
G(s):=\int\frac{\mu(dv)}{v^2+s}=\frac1{\sqrt s}\,\frac{\xi'}{\xi}(\theta'+\sqrt s).
\]
On a compact subset of $\cut$ one has $|v^2+s|\ge c(1+v^2)$, so $G$ is holomorphic on $\cut$. The right side, with the principal square root, is meromorphic on $\cut$, and the identity theorem on the connected set $\cut$ minus its poles makes $G$ equal to it there. Squaring maps $\Re w>0$ onto $\cut$, so $s_0=(\rho-\theta')^2\in\cut$ with $\sqrt{s_0}=\rho-\theta'$, and near $s_0$ the right side is $2m_\rho/(s-s_0)+O(1)$, where $m_\rho\ge1$ is the multiplicity of $\rho$, while $G$ is bounded near $s_0$. This contradiction proves (ii).
\end{proof}

The theorem is the Thorin criterion of Chapter~\ref{chap:lv} in dynamical form, with part (ii) proved by the pole argument of Theorem~\ref{lv:thm:thorin}: RH is the convergence of the backward flow from the explicit level one measure at every level $\theta'>\tfrac12$. The flow relocates the difficulty into the multiplier $e^{(1-\theta')|\omega|}$, which amplifies whatever is not known about $\widehat{\mathfrak U_1}$ at high frequency. On the prime side, high frequency means large $n$: the content of $\widehat{\mathfrak U_1}$ on $|\omega|\le\Omega$ is $\sum_{n\le e^\Omega}\Lambda(n)/n$, and the backward flow to level $\theta'$ heuristically asks for $\sum_{n\le x}\Lambda(n)n^{-\theta'}$ to behave, which is the statement $\psi(x)-x=O(x^{\theta'+\varepsilon})$ for every $\varepsilon>0$ (the classical equivalence with the zero free half plane $\Re s>\theta'$).

\subsection{Numerical backward flow}\label{lo:sec:numflow}

We sampled $u_1$ on $[0,400]$ at spacing $0.05$ from \eqref{lo:eq:dens} with $\theta=1$, evaluating $\zeta$ and $\zeta'$ on the line $\Re s=1$ only (\texttt{python-flint}, $128$ bits), with the analytic value $u_1(0)$ at $v=0$, extended by symmetry, applied a Hann window and the multiplier of \eqref{lo:eq:backward} with a sharp cutoff at $\Omega$, not the Fej\'er window of \eqref{lo:eq:backward}, inverted and divided by the window. A peak is a local maximum of the output in $100<v<300$ of height at least a quarter of the maximum, and the peaks are compared with the $109$ zeros of $\zeta$ in that range ($N(300)-N(100)=138-29$). The Fej\'er variant multiplies the sharp multiplier by $1-|\omega|/\Omega$.

\begin{table}[ht]
\centering
\small
\begin{tabular}{ccccccc}
\toprule
$\theta'$ & $\Omega$ & cutoff & peaks & zeros & max error & min value\\
\midrule
0.75 & 40 & sharp & 109 & 109 & 0.035 & $+0.075$\\
0.50 & 40 & sharp & 109 & 109 & 0.027 & $-3.2$\\
0.50 & 60 & sharp & 213 & 109 & 0.042 & $-16.5$\\
0.50 & 60 & Fej\'er & 109 & 109 & 0.025 & $-0.020$\\
\bottomrule
\end{tabular}
\caption{The Fourier backward flow from samples of $u_1$ on $\Re s=1$ (computed). The error is the largest distance between a zero in $100<v<300$ and the nearest peak. At $\Omega=60$ the sharp cutoff has $213$ peaks of height at least a quarter of the maximum, $109$ of them at the zeros, and the maximum is $+26.6$.}\label{lo:tab:flow}
\end{table}

\begin{resultpar}\label{lo:res:flow}At $\theta'=0.75$ the flow produces a positive function with $109$ peaks, each within $0.035$ of a zero, as Theorem~\ref{lo:thm:dichotomy}(i) predicts if $\zeta$ has no zeros in $\Re s>3/4$: the level $\tfrac34$ measure is then a smooth positive density whose peaks lie near the zeros. At $\theta'=0.5$ the peaks lie within $0.027$ of the zeros, but the function oscillates below zero between them, because of the sharp cutoff: its kernel $\sin(\Omega v)/(\pi v)$ has negative lobes, the first of depth about $0.22\,\Omega/\pi\approx2.8$ at $\Omega=40$, whereas under RH the Fej\'er approximant $\mathfrak U_{1/2}*\Phi_\Omega$ of Theorem~\ref{lo:thm:dichotomy}(i) is nonnegative, because $\Phi_\Omega\ge0$. Theorem~\ref{lo:thm:dichotomy} does not cover the sharp cutoff. The Fej\'er run at $\Omega=60$ is clean: $109$ peaks, error $0.025$ and minimum $-0.020$ (the small negative values come from the Hann window and its division); at $\Omega=40$ its minimum is $-0.011$.\end{resultpar}

The cutoff is limited by the accuracy of the samples: at level $\tfrac12$ the multiplier at the cutoff is $e^{\Omega/2}$, about $10^{13}$ at $\Omega=60$, so an error of $10^{-7}$ in the single sample at $v=0$ already produces oscillations of order $10^2$ in $100<v<300$. This is what the source printed: its $\Omega=60$ row, with $1022$ peaks, error $1.9$ and minimum $-430$, has its minimum reproduced only if the sample at $v=0$, the pole of $\zeta$, is replaced by the sample at $v=0.05$ (an error of $8.8\times10^{-8}$; minimum $-430.2$, maximum $+430.9$; the peak count and the error are not reproduced, the source giving no peak rule), and the source read it as amplified rounding error. With the analytic $u_1(0)$ the output of the sharp cutoff at $\Omega=60$ is not noise but the ringing of the cutoff: each atom gives a sinc of height $\Omega/\pi\approx19$ and the sidelobes of $213$ peaks follow; double precision rounding alone would give oscillations of order $3\times10^{-4}$. The line $\Re s=1$ determines the zeros, and the $109$ zeros tested are read off it to $0.025$ to $0.035$ (two decimals, not the three decimals and the first two hundred zeros of the source), but at level $\tfrac12$ the precision required of the level one data grows like $e^{\Omega/2}$ with the cutoff.
